\documentclass[10pt,a4paper]{amsart}
\usepackage{amsmath,amssymb,amsthm,graphicx}
% Retained preamble for the Stacks Project chapter {CH} (shared project preamble).
% The chapter file pinned-original-source.tex only carries \input{{preamble}} on its first line,
% so no preamble text lives inside the chapter itself.  The Stacks Project uses one single shared
% preamble for every chapter: the snapshot stages it three times (sources/stacks-chow/preamble.tex,
% sources/stacks-etale-cohomology/preamble.tex and focused-seeds/AGP/sources/stacks-snapshot-20261011/
% preamble.tex) and those copies are identical in size.  The block below is copied mechanically,
% line for line, from sources/stacks-chow/preamble.tex lines 14-18 and 154-217: the Xy-pic setup and
% the native theorem declarations together with the mathematical macros.  Nothing was added, removed
% or redefined; the wrapper class and the amsmath/amssymb/amsthm/graphicx base packages are supplied
% by the assembler.
% For commutative diagrams we use Xy-pic
% Chapter: stacks-spaces-more-morphisms.tex -- More on Morphisms of Spaces.
\usepackage[all]{xy}

% We use 2cell for 2-commutative diagrams.
\xyoption{2cell}

% Theorem environments.
%
\theoremstyle{plain}
\newtheorem{theorem}[subsection]{Theorem}
\newtheorem{proposition}[subsection]{Proposition}
\newtheorem{lemma}[subsection]{Lemma}

\theoremstyle{definition}
\newtheorem{definition}[subsection]{Definition}
\newtheorem{example}[subsection]{Example}
\newtheorem{exercise}[subsection]{Exercise}
\newtheorem{situation}[subsection]{Situation}

\theoremstyle{remark}
\newtheorem{remark}[subsection]{Remark}
\newtheorem{remarks}[subsection]{Remarks}

\numberwithin{equation}{subsection}

% Macros
%
\def\lim{\mathop{\mathrm{lim}}\nolimits}
\def\colim{\mathop{\mathrm{colim}}\nolimits}
\def\Spec{\mathop{\mathrm{Spec}}}
\def\Hom{\mathop{\mathrm{Hom}}\nolimits}
\def\Ext{\mathop{\mathrm{Ext}}\nolimits}
\def\SheafHom{\mathop{\mathcal{H}\!\mathit{om}}\nolimits}
\def\SheafExt{\mathop{\mathcal{E}\!\mathit{xt}}\nolimits}
\def\Sch{\mathit{Sch}}
\def\Mor{\mathop{\mathrm{Mor}}\nolimits}
\def\Ob{\mathop{\mathrm{Ob}}\nolimits}
\def\Sh{\mathop{\mathit{Sh}}\nolimits}
\def\NL{\mathop{N\!L}\nolimits}
\def\CH{\mathop{\mathrm{CH}}\nolimits}
\def\proetale{{pro\text{-}\acute{e}tale}}
\def\etale{{\acute{e}tale}}
\def\QCoh{\mathit{QCoh}}
\def\Ker{\mathop{\mathrm{Ker}}}
\def\Im{\mathop{\mathrm{Im}}}
\def\Coker{\mathop{\mathrm{Coker}}}
\def\Coim{\mathop{\mathrm{Coim}}}

% Boxtimes
%
\DeclareMathSymbol{\boxtimes}{\mathbin}{AMSa}{"02}

%
% Macros for moduli stacks/spaces
%
\def\QCohstack{\mathcal{QC}\!\mathit{oh}}
\def\Cohstack{\mathcal{C}\!\mathit{oh}}
\def\Spacesstack{\mathcal{S}\!\mathit{paces}}
\def\Quotfunctor{\mathrm{Quot}}
\def\Hilbfunctor{\mathrm{Hilb}}
\def\Curvesstack{\mathcal{C}\!\mathit{urves}}
\def\Polarizedstack{\mathcal{P}\!\mathit{olarized}}
\def\Complexesstack{\mathcal{C}\!\mathit{omplexes}}
% \Pic is the operator that assigns to X its picard group, usage \Pic(X)
% \Picardstack_{X/B} denotes the Picard stack of X over B
% \Picardfunctor_{X/B} denotes the Picard functor of X over B
\def\Pic{\mathop{\mathrm{Pic}}\nolimits}
\def\Picardstack{\mathcal{P}\!\mathit{ic}}
\def\Picardfunctor{\mathrm{Pic}}
\def\Deformationcategory{\mathcal{D}\!\mathit{ef}}
\usepackage{cleveref}
\crefname{originallemma}{original lemma}{original lemmas}
\Crefname{originallemma}{Original lemma}{Original lemmas}
\makeatletter
\newlabel{lemma-thicken-property-morphisms-cartesian}{{09ZZ}{source}{}{sourcepointer-1}{}}
\newlabel{lemma-thicken-property-morphisms-cartesian@cref}{{[originallemma][0][]09ZZ}{[1][0][]source}{}{}{}}
\makeatother
\crefname{originalremark}{original remark}{original remarks}
\Crefname{originalremark}{Original remark}{Original remarks}
\makeatletter
\newlabel{remark-reformulate-existence-theorem}{{08BF}{source}{}{sourcepointer-2}{}}
\newlabel{remark-reformulate-existence-theorem@cref}{{[originalremark][0][]08BF}{[1][0][]source}{}{}{}}
\makeatother
\crefname{originaltheorem}{original theorem}{original theorems}
\Crefname{originaltheorem}{Original theorem}{Original theorems}
\makeatletter
\newlabel{theorem-grothendieck-existence}{{08BE}{source}{}{sourcepointer-3}{}}
\newlabel{theorem-grothendieck-existence@cref}{{[originaltheorem][0][]08BE}{[1][0][]source}{}{}{}}
\makeatother
\title{Algebraization of a formal finite and proper algebraic space over a complete Noetherian base}
\author{The Stacks Project Authors}
\date{stacks-project:spaces-more-morphisms.tex:local-snapshot-20261011; GNU FDL 1.2 or later}
\begin{document}
\maketitle
\noindent\textit{Editorial note.} Counted span is the author's own proof environment, source lines 10216-10236, immediately following the statement at 10193-10214.\par\medskip
\noindent\textit{Editorial note.} The proof uses Grothendieck's existence theorem twice: once in the reformulated form to produce the module, once as full faithfulness to algebraize it; both are navigation pointers into the same chapter.\par\medskip
\noindent\textit{Editorial note.} The final step setting Y equal to the relative spectrum is displayed and completed inside the span.\par\medskip
\noindent\textit{Allowed standard prerequisites (editorial).}\begin{itemize}
\item Thickening property for cartesian squares of morphisms. chapter lemma lemma-thicken-property-morphisms-cartesian, cited at proof line 10220. Shows each vertical map f\_n : Y\_n -> X\_n is finite because its base change to X\_1 equals the given finite f\_1.
\item Reformulation of Grothendieck's existence theorem. chapter remark remark-reformulate-existence-theorem, cited at proof line 10226. Produces the coherent O\_X-module F restricting to the given F\_n on each X\_n and records that its support is proper over S.
\item Grothendieck's existence theorem: the completion functor is fully faithful. chapter theorem theorem-grothendieck-existence, cited at proof line 10231. Lets the compatible multiplication maps on the F\_n descend to an algebra structure on F.
\item Relative spectrum of a sheaf of algebras. used at proof line 10235. Realises the algebraized object as Y = Spec\_X(F), which is finite over X and proper over S.
\end{itemize}\medskip
\ifdefined\hypertarget\hypertarget{sourcepointer-1}{}\fi
\noindent\textit{Source reference.} Original lemma 09ZZ: sources/stacks-spaces-more-morphisms/pinned-original-source.tex, source lines 2044--2044. Navigation pointer only: the counted proof cites this statement by \textbackslash{}ref as a standing input; the deduction the counted proof has to make is carried out inside the counted span, and no part of that deduction is delegated to this pointer.\par\medskip
\ifdefined\hypertarget\hypertarget{sourcepointer-2}{}\fi
\noindent\textit{Source reference.} Original remark 08BF: sources/stacks-spaces-more-morphisms/pinned-original-source.tex, source lines 10085--10085. Navigation pointer only: the counted proof cites this statement by \textbackslash{}ref as a standing input; the deduction the counted proof has to make is carried out inside the counted span, and no part of that deduction is delegated to this pointer.\par\medskip
\ifdefined\hypertarget\hypertarget{sourcepointer-3}{}\fi
\noindent\textit{Source reference.} Original theorem 08BE: sources/stacks-spaces-more-morphisms/pinned-original-source.tex, source lines 9990--9990. Navigation pointer only: the counted proof cites this statement by \textbackslash{}ref as a standing input; the deduction the counted proof has to make is carried out inside the counted span, and no part of that deduction is delegated to this pointer.\par\medskip
\begin{lemma}
\label{lemma-algebraize-formal-algebraic-space-finite-over-proper}
Let $A$ be a Noetherian ring complete with respect to an ideal $I$.
Write $S = \Spec(A)$ and $S_n = \Spec(A/I^n)$.
Let $X \to S$ be a morphism of algebraic spaces that is separated
and of finite type.
For $n \geq 1$ we set $X_n = X \times_S S_n$.
Suppose given a commutative diagram
$$
\xymatrix{
Y_1 \ar[r] \ar[d] & Y_2 \ar[r] \ar[d] & Y_3 \ar[r] \ar[d] & \ldots \\
X_1 \ar[r]^{i_1} & X_2 \ar[r]^{i_2} & X_3 \ar[r] & \ldots
}
$$
of algebraic spaces with cartesian squares. Assume that
\begin{enumerate}
\item $Y_1 \to X_1$ is a finite morphism, and
\item $Y_1 \to S_1$ is proper.
\end{enumerate}
Then there exists a finite morphism of algebraic spaces $Y \to X$ such that
$Y_n = Y \times_S S_n$ for all $n \geq 1$. Moreover, $Y$ is proper over $S$.
\end{lemma}

\begin{proof}
Let's write $f_n : Y_n \to X_n$ for the vertical morphisms.
As the squares in the statement are cartesian
we see that the base change of $f_n$ to $X_1$ is $f_1$.
Thus Lemma \ref{lemma-thicken-property-morphisms-cartesian}
shows that $f_n$ is a finite morphism.
Set $\mathcal{F}_n = f_{n, *}\mathcal{O}_{Y_n}$.
Using that the squares are cartesian we see that
the pullback of $\mathcal{F}_{n + 1}$ to $X_n$ is $\mathcal{F}_n$.
Hence Grothendieck's existence theorem, as reformulated in
Remark \ref{remark-reformulate-existence-theorem},
tells us there exists a coherent $\mathcal{O}_X$-module $\mathcal{F}$
whose restriction to $X_n$ recovers $\mathcal{F}_n$.
Moreover, the support of $\mathcal{F}$ is proper over $S$.
As the completion functor is fully faithful
(Theorem \ref{theorem-grothendieck-existence})
we see that the multiplication maps
$\mathcal{F}_n \otimes_{\mathcal{O}_{X_n}} \mathcal{F}_n \to
\mathcal{F}_n$ fit together to give an algebra structure on $\mathcal{F}$.
Setting $Y = \underline{\Spec}_X(\mathcal{F})$ finishes the proof.
\end{proof}

\end{document}
